\chapter*{Bibliografía}
\phantomsection
\addcontentsline{toc}{chapter}{Bibliografía}
\label{cap:bibliografia}

Durante el desarrollo del proyecto se consultaron los siguientes libros
para lograr una mejor comprensión de los fundamentos y las técnicas de
construcción de compiladores. En particular, \emph{Compiladores:
principios, técnicas y herramientas}, de Aho y colaboradores,
\emph{Crafting a Compiler} y \emph{Engineering a Compiler} fueron las
referencias principales para acompañar el estudio de las distintas
etapas y orientar las decisiones de implementación. Los demás libros
sirvieron como material de consulta complementario.

\begin{enumerate}[leftmargin=*,label={[\arabic*]}]
  \item Aho, Alfred V.; Lam, Monica S.; Sethi, Ravi; y Ullman, Jeffrey D.
  \emph{Compiladores: principios, técnicas y herramientas}.
  Segunda edición. Pearson Educación, 2008.

  \item Fischer, Charles N.; Cytron, Ron K.; y LeBlanc, Richard J., Jr.
  \emph{Crafting a Compiler}. Addison-Wesley, 2010.

  \item Cooper, Keith D.; y Torczon, Linda.
  \emph{Engineering a Compiler}. Segunda edición. Morgan Kaufmann, 2012.

  \item Appel, Andrew W., con Palsberg, Jens.
  \emph{Modern Compiler Implementation in Java}.
  Segunda edición. Cambridge University Press, 2002.

  \item Watt, David A.; y Brown, Deryck F.
  \emph{Programming Language Processors in Java: Compilers and Interpreters}.
  Prentice Hall, 2000.

  \item Grune, Dick; van Reeuwijk, Kees; Bal, Henri E.; Jacobs, Ceriel J. H.;
  y Langendoen, Koen. \emph{Modern Compiler Design}.
  Segunda edición. Springer, 2012.
\end{enumerate}
